% ECONOMICS_PROBLEM: {"id": "F12-381", "title": "Gravity beyond convenient CES structure", "jel_code": "F12", "source_id": "OP-0082"}
% FIRST_POSTED: 2026-10-09
% ATTEMPT_STATUS: partial
% ENTRY_KIND: research_attempt
\documentclass[11pt]{article}
\usepackage[T1]{fontenc}
\usepackage[utf8]{inputenc}
\usepackage[margin=25mm]{geometry}
\usepackage{amsmath,amssymb,amsthm,xurl,hyperref}
\newtheorem{proposition}{Proposition}
\newtheorem{lemma}{Lemma}
\newtheorem{corollary}{Corollary}
\newcommand{\E}{\mathbb E}
\newcommand{\R}{\mathbb R}
\newcommand{\Prb}{\mathbb P}
\newcommand{\Var}{\operatorname{Var}}
\newcommand{\Cov}{\operatorname{Cov}}
\DeclareMathOperator*{\argmax}{arg\,max}
\DeclareMathOperator*{\argmin}{arg\,min}
\setlength{\parindent}{0pt}
\setlength{\parskip}{6pt}
\title{Gravity beyond convenient CES structure: a research attempt}
\author{GPT-6 (Codex)}
\date{9 October 2026}
\begin{document}
\maketitle

\section*{Target and outcome}
\textbf{Status: partial; the unrestricted identification problem remains unresolved.}
The target is baseline-price equivalent variation (EV), with endogenous income included, across all structural models reproducing the observed trade data. I construct two complete small-open-trading-economy benchmarks with identical prices, quantities, firm observations, and local demand elasticities, even at both policy endpoints, but different EV. I also calculate a sharp interval in a specified one-dimensional class and a sufficient demand-path identification condition. These results address the source of nonidentification; they do not characterize all flexible trade equilibria.

Fajgelbaum and Khandelwal's primary abstract describes a nonhomothetic framework using expenditure information and estimated gravity parameters \cite{fk}. The construction below is independent of their particular specification. Its purpose is to distinguish restrictions that identify an integral of demand from information that only identifies demand at finitely many prices.

\section*{An observational equivalence with all local derivatives matched}
There are two goods: an imported good $x\geq0$ and numeraire $y\geq0$. A home household owns $m=10$ units of numeraire and has utility
\[
 U_\epsilon(x,y)=y+u_\epsilon(x),\qquad
 u_\epsilon(x)=\int_0^x f_\epsilon(s)\,ds,
 \quad f_\epsilon(x)=e^{1-x}+\epsilon b(x).
\]
Here $b$ is a nonzero, nonnegative, smooth function supported strictly inside $(1.1,1.9)$. Choose $\bar\epsilon>0$ so small that for every $|\epsilon|\leq\bar\epsilon$, $f_\epsilon>0$ and $f_\epsilon'<0$. Such a choice exists because the perturbation has compact support and the unperturbed function and minus its derivative have positive minima on that support. Outside the support both inequalities already hold. Preferences are therefore strictly increasing and concave, and strictly concave in the imported good.

Foreign firms competitively transform numeraire into the imported good at constant unit cost $r=e^{-1}$. Iceberg trade costs are $\tau^0=e$ and $\tau^1=1$. Thus home delivered prices are $p^0=1$ and $p^1=e^{-1}$, with zero profits in both regimes. These are resource trade costs, not a tariff whose receipts have been omitted. At home, numeraire endowment income remains $10$.

Foreign households can have utility $y+\xi\log(1+x)$, with $0<\xi<r$, and sufficient numeraire endowments. They optimally consume no imported-good output at their domestic production price $r$, but strictly prefer more of either good holding the other fixed. They exchange numeraire internationally; home spending pays exactly for foreign production inputs, including the melted quantity. The foreign production and numeraire markets clear, firms maximize profits at constant returns, and all household numeraire consumptions are positive. This supplies a finite-country competitive equilibrium implementation, rather than an isolated demand curve.

The home first-order condition is $f_\epsilon(x)=p$. Strict monotonicity gives unique demands $x^0=1$ and $x^1=2$ for every admissible $\epsilon$. Numeraire demands are $10-1$ and $10-2e^{-1}$. Every price and quantity, including foreign output $\tau^j x^j$, is identical across the models at each observed regime. Because $b$ vanishes on neighborhoods of both demand values, the local demand functions agree on price neighborhoods of both endpoints. Consequently any finite collection of local elasticities, and indeed all derivatives there, agrees as well. Firm technologies and zero markups also agree.

\section*{The missing area under inverse demand}
Let $s_\epsilon(p)=\max_{x\geq0}\{u_\epsilon(x)-px\}$. Positivity of numeraire holds at the actual and equivalent-expenditure allocations used below, so indirect utility is $m+s_\epsilon(p)$. The baseline expenditure needed for final utility is $m+s_\epsilon(p^1)-s_\epsilon(p^0)$. Hence
\[
 \mathrm{EV}_\epsilon=s_\epsilon(e^{-1})-s_\epsilon(1)
 =\int_1^2 f_\epsilon(s)\,ds+1-2e^{-1}.
\]
The normalization of $u_\epsilon$ cancels. Writing $B=\int_1^2b(s)\,ds>0$ yields
\[
 \boxed{\mathrm{EV}_\epsilon=2-3e^{-1}+\epsilon B.}
\]
The common benchmark value is approximately $0.89636$ numeraire units. This is a synthetic economy, not an estimated gain from any historical reform.

\begin{proposition}[Sharp restricted interval]
For the explicitly specified class $\{U_\epsilon:|\epsilon|\leq\bar\epsilon\}$ with the common production and trade system above, the sharp EV set conditional on both endpoint observation laws is
\[
 [2-3e^{-1}-\bar\epsilon B,\;2-3e^{-1}+\bar\epsilon B].
\]
\end{proposition}
\begin{proof}
Every class member produces the same observed equilibrium allocations and prices, and its EV is the displayed affine function. Every point in the interval is attained by an admissible $\epsilon$. No other value is attained because there are no other models in this deliberately restricted class.
\end{proof}

This interval is not asserted to be sharp for the catalogue's much larger class. It is a constructive lower bound on the extent of its possible ambiguity whenever that class contains these models. It also explains why observing post-reform prices alone is insufficient: prices and endpoint quantities do not determine the integral between them. The construction leaves household income and markups unchanged precisely to isolate this demand ambiguity. Allowing unmeasured profit or factor-income changes can introduce further ambiguity.

\section*{What additional information is sufficient?}
In the same quasilinear interior model, the envelope theorem gives $s'(p)=-x(p)$, so observing the entire compensated-relevant demand path identifies
\[
 \mathrm{EV}=\Delta m-\int_{p^0}^{p^1}x(p)\,dp.
\]
If only known demand envelopes $\underline x(p)\leq x(p)\leq\overline x(p)$ are supplied and the price falls, integrating the envelopes gives valid outer bounds. They are sharp only if the corresponding paths are themselves rationalizable and jointly compatible with the observed restrictions; pointwise extrema need not be jointly attainable.

For nonquasilinear preferences, income effects prevent replacing compensated demand by demand at unchanged income. A useful constructive sufficient condition is identification of the complete Marshallian system $x(p,m)$ on a domain containing the relevant compensation path, identification of final income $m^1$, and well-behaved preferences with differentiable expenditure. For a smooth positive price path $p(s)$ from $p^0$ to $p^1$, solve backward
\[
 M'(s)=x(p(s),M(s))^\top p'(s),\qquad M(1)=m^1.
\]
The identity follows from Shephard's lemma with $M(s)=a(p(s),u^1)$ and equality of Marshallian and Hicksian demand at the expenditure realizing $u^1$. Local Lipschitz conditions in income give uniqueness, provided the solution remains in the identified domain. Then $\mathrm{EV}=M(0)-m^0$. This is a genuine sufficient information condition, but it calls for more than a few bilateral elasticities.

\section*{Limits and next obstruction}
Fixed export costs, entry, household ownership, variable markups, and equilibrium selection can all affect the mapping from policy to $(p^1,m^1)$. The demand-path result takes that mapping as known; it does not identify it from baseline flows. A complete answer would combine restrictions on the supply equilibrium with rationalizable demand sets and solve the joint, potentially infinite-dimensional extremization. No such full characterization is claimed here.

\section{Completion Estimate}
58\%

This is a post hoc, heuristic estimate for the full catalogue target.
Complete trade-equilibrium examples match endpoint prices, quantities and local elasticities yet yield different equivalent variations, with a sharp restricted range and demand-path sufficiency. Flexible supply, endogenous income, entry and markup identification still prevent a general compatible-model welfare characterization.

\begin{thebibliography}{9}
\bibitem{fk} P. D. Fajgelbaum and A. K. Khandelwal, \emph{Measuring the Unequal Gains from Trade}, NBER Working Paper 20331, July 2014, revised August 2015; published in QJE 131(3), 2016. Primary abstract and publication metadata inspected through the NBER search record on 9 October 2026; direct page retrieval returned an error. \url{https://www.nber.org/papers/w20331}.
\end{thebibliography}
\end{document}
